% ===========================================================================
\chapter{Modelagem Estrutural do Frontend (Angular 18+ e Signals)}
\label{cap_modelagem_frontend}
% ===========================================================================

A camada de interface do usuário do SIGEA-GTT é concebida sobre o ecossistema moderno do Angular 18+, adotando o paradigma de *Standalone Components* e a reatividade granular de *Angular Signals*. Ao contrário das arquiteturas web convencionais centradas em módulos imperativos pesados (\texttt{NgModule}), o frontend do SIGEA-GTT opera com componentes desacoplados, serviços injetáveis no escopo raiz (\texttt{providedIn: 'root'}) e controle reativo de estado sem sobrecarga de inscrições manuais em observáveis.

% ===========================================================================
\section{Diagrama de Classes e Componentes do Frontend}
\label{sec_diagrama_classes_frontend}
% ===========================================================================

A Figura \ref{fig_diagrama_classes_frontend} exibe a modelagem estrutural das principais classes, serviços e guardas de rota do cliente Angular:

\begin{figure}[!ht]
\centering
\caption{Diagrama de Classes e Componentes do Frontend (Angular 18+)}
\label{fig_diagrama_classes_frontend}
\resizebox{0.96\textwidth}{!}{%
\begin{tikzpicture}[
    service_box/.style={
        rectangle split,
        rectangle split parts=3,
        rectangle split part fill={clinical-100!30, white, white},
        draw=clinical-700,
        thick,
        rounded corners=2pt,
        font=\tiny,
        minimum width=5.2cm,
        text width=5.0cm,
        align=left
    },
    comp_box/.style={
        rectangle split,
        rectangle split parts=3,
        rectangle split part fill={teal!15, white, white},
        draw=teal!80!black,
        thick,
        rounded corners=2pt,
        font=\tiny,
        minimum width=4.8cm,
        text width=4.6cm,
        align=left
    },
    guard_box/.style={
        rectangle split,
        rectangle split parts=3,
        rectangle split part fill={alert-amber!15, white, white},
        draw=alert-amber!90!black,
        thick,
        rounded corners=2pt,
        font=\tiny,
        minimum width=4.8cm,
        text width=4.6cm,
        align=left
    },
    dep/.style={
        draw=clinical-900,
        thick,
        dashed,
        -Latex
    }
]

    % ====================================================
    % SERVIÇO DE AUTENTICAÇÃO E SESSÃO
    % ====================================================
    \node[service_box] (auth_serv_front) at (-4.0, 5.0) {
        \centering\textbf{\scriptsize <<Injectable>>\\AuthService}
        \nodepart{two}
        - currentUser : WritableSignal<User | null>\\
        + isAuthenticated : Signal<boolean>\\
        + userRole : Signal<UserRole | null>\\
        + mustChangePassword : Signal<boolean>
        \nodepart{three}
        + login(credentials) : Observable<LoginResponse>\\
        + logout() : void\\
        + getToken() : string | null\\
        + clearSession() : void
    };

    % ====================================================
    % SERVIÇO DE AUDITORIA CLÍNICA
    % ====================================================
    \node[service_box] (audit_serv_front) at (4.0, 5.0) {
        \centering\textbf{\scriptsize <<Injectable>>\\AuditSessionService}
        \nodepart{two}
        + remainingSeconds : WritableSignal<number>\\
        + isPaused : WritableSignal<boolean>\\
        + activeTriggers : WritableSignal<Trigger[]>\\
        + qualityTools : WritableSignal<QualityTools>
        \nodepart{three}
        + startTimer() : void\\
        + pauseTimer() : void\\
        + resumeTimer() : void\\
        + addTrigger(Trigger) : void\\
        + submitAudit() : Observable<void>
    };

    % ====================================================
    % COMPONENTES DE TELA
    % ====================================================
    \node[comp_box] (shell_comp) at (-4.0, 0.0) {
        \centering\textbf{\scriptsize <<Component>>\\AppShellComponent}
        \nodepart{two}
        - authService : AuthService\\
        - sidebarOpen : WritableSignal<boolean>
        \nodepart{three}
        + toggleSidebar() : void\\
        + handleLogout() : void
    };

    \node[comp_box] (audit_comp) at (4.0, 0.0) {
        \centering\textbf{\scriptsize <<Component>>\\AuditWorkspaceComponent}
        \nodepart{two}
        - auditService : AuditSessionService\\
        - timerColor : Signal<string>
        \nodepart{three}
        + onTriggerSelected(Trigger) : void\\
        + onConfirmSubmit() : void
    };

    % ====================================================
    % GUARDAS E INTERCEPTORES
    % ====================================================
    \node[guard_box] (auth_guard) at (-4.0, -4.5) {
        \centering\textbf{\scriptsize <<Guard>>\\AuthGuard / RoleGuard}
        \nodepart{two}
        - authService : AuthService\\
        - router : Router
        \nodepart{three}
        + canActivate(route, state) : boolean
    };

    \node[guard_box] (auth_interceptor) at (4.0, -4.5) {
        \centering\textbf{\scriptsize <<Interceptor>>\\AuthInterceptor}
        \nodepart{two}
        - authService : AuthService
        \nodepart{three}
        + intercept(req, next) : Observable<HttpEvent>
    };

    % Dependências
    \draw[dep] (shell_comp) -- (auth_serv_front);
    \draw[dep] (audit_comp) -- (audit_serv_front);
    \draw[dep] (auth_guard) -- (auth_serv_front);
    \draw[dep] (auth_interceptor) -- (auth_serv_front);

\end{tikzpicture}
}
\fonte{Elaborado pelo autor (2026).}
\end{figure}

% ===========================================================================
\section{Gerenciamento Reativo de Estado com Angular Signals}
\label{sec_angular_signals}
% ===========================================================================

No SIGEA-GTT, o estado da aplicação no cliente é governado primariamente por primitivas reativas de **Signals**:
\begin{itemize}
    \item \texttt{WritableSignal<T>}: Armazena valores reativos mutáveis (ex.: \texttt{remainingSeconds} do cronômetro de 20 minutos, lista de gatilhos adicionados);
    \item \texttt{Signal<T> (Computed)}: Deriva valores computados automaticamente a partir de outros sinais com avaliação preguiçosa (\textit{lazy evaluation}) e memoização, tal como o sinal derivado \texttt{isAuthenticated}, que recalcula seu estado estritamente quando o sinal base \texttt{currentUser} for alterado;
    \item \texttt{Effect}: Dispara efeitos colaterais reativos em resposta a mudanças de estado, como a alteração das cores do cronômetro de auditoria para amarelo (\texttt{\#D97706}) aos 15 minutos e carmim cirúrgico (\texttt{\#DC2626}) aos 18 minutos.
\end{itemize}

% ===========================================================================
\section{Segurança de Rotas e Interceptores HTTP}
\label{sec_seguranca_rotas_frontend}
% ===========================================================================

A segurança da interface é garantida pela atuação conjunta de Guardas funcionais e Interceptores:
\begin{itemize}
    \item \textbf{\texttt{AuthInterceptor}}: Injeta o cabeçalho \texttt{Authorization:}\linebreak[0] \texttt{Bearer <token>} em todas as requisições direcionadas à API do backend, capturando respostas HTTP 401 (\textit{Unauthorized}) para expirar a sessão e redirecionar ao login;
    \item \textbf{\texttt{AuthGuard} e \texttt{RoleGuard}}: Inspecionam as rotas no momento da navegação. Caso o usuário não esteja autenticado ou não possua o perfil exigido na rota (\texttt{ADMIN}, \texttt{PROFESSOR} ou \texttt{STUDENT}), o acesso é sumariamente bloqueado;
    \item \textbf{\texttt{FirstLoginGuard}}: Verifica se a flag \texttt{mustChangePassword} é verdadeira, forçando o redirecionamento imediato para a tela de troca de senha provisória caso o discente ou docente tente acessar outras áreas da plataforma.
\end{itemize}
